% frontmatter/english/palabras.tex -- para la niña: las palabras de las
% instrucciones, cada una con el icono que lleva su actividad en el
% título de la caja (preamble-english.tex), y las palabras de lugar de
% "Where is it?" (in, on, under...) dibujadas. Quien empieza a leer en
% inglés reconoce antes el dibujo que la palabra; esta página es el
% puente entre los dos, y se lee con el adulto el primer día.
\section*{The words in this book}

{\large Every activity has a picture. Read the words with a grown-up.\par}
\vspace{5mm}

% El icono va en un \makebox de ancho fijo: algunos (rodea, ordena) son
% más anchos que una letra, y así sobresalen un poco a los lados en vez
% de empujar la columna.
\newcommand{\palabraLibro}[4]{%
  \makebox[17mm][c]{\fontsize{26}{26}\selectfont\color{#1}#2}\hspace{2mm}%
  \parbox[c]{0.33\linewidth}{\raggedright{\Large\bfseries #3}\\[0.4mm]{\large #4}}}
{\renewcommand{\arraystretch}{1.4}
\noindent\begin{tabular}{@{}l@{\hspace{6mm}}l@{}}
\palabraLibro{colorLectura}{\icoLibro}{read}{Read the text.} &
\palabraLibro{colorDibuja}{\icoLapiz}{draw}{Draw a cat.} \\[3mm]
\palabraLibro{colorCompleta}{\icoCera}{colour}{Colour the cat black.} &
\palabraLibro{colorDibuja}{\icoSitio}{Where is it?}{Draw it in the right place.} \\[3mm]
\palabraLibro{colorResponde}{\icoMarca}{tick}{Tick yes or no.} &
\palabraLibro{colorResponde}{\icoRodea}{circle}{Circle the right answer.} \\[3mm]
\palabraLibro{colorRelaciona}{\icoUne}{match}{Draw a line to join them.} &
\palabraLibro{colorResponde}{\icoOrden}{What happened first?}{Write 1, 2, 3.} \\[3mm]
\palabraLibro{colorResponde}{\icoHueco}{the missing word}{Write the word that is missing.} &
\palabraLibro{colorResponde}{\icoBusca}{find}{Find the words in the text.} \\[3mm]
\palabraLibro{colorRelaciona}{\icoAdivina}{riddle}{What is it? Guess!} &
\palabraLibro{colorDibuja}{\icoVinetas}{the story}{Draw the beginning, the middle and the end.} \\[3mm]
\palabraLibro{colorCrea}{\icoEstrella}{look back}{Can you remember?} &
\palabraLibro{colorCrea}{\icoCara}{your turn}{Now it's you!} \\
\end{tabular}\par}

\vspace{9mm}
{\Large\bfseries Where is the ball?\par}
\vspace{3mm}

% Una caja, una mesa y una pelota, en cada una de las posiciones que
% salen en "Where is it?". Detrás = la caja tapa media pelota; delante =
% la pelota tapa la caja.
\newcommand{\cajaLugar}[1]{\filldraw[fill=white] (#1) rectangle ++(1.4,1.1);
  \draw ($(#1)+(0,1.1)$) -- ++(0.35,0.28) -- ++(1.4,0) -- ++(-0.35,-0.28);
  \draw ($(#1)+(1.4,0)$) -- ++(0.35,0.28) -- ++(0,1.1);}
\newcommand{\pelotaLugar}[1]{\filldraw[fill=colorCompletaClaro] (#1) circle (0.36);
  \draw ($(#1)+(-0.34,0.1)$) .. controls ($(#1)+(0,-0.12)$) .. ($(#1)+(0.34,0.1)$);}
\newcommand{\mesaLugar}[1]{\draw (#1) ++(0,1.3) -- ++(2.2,0);
  \draw (#1) ++(0.2,1.3) -- ++(0,-1.3); \draw (#1) ++(2,1.3) -- ++(0,-1.3);}
\begin{center}
\begin{tikzpicture}[line width=1.4pt, line cap=round, line join=round, scale=1.35,
  etiqueta/.style={font=\Large\bfseries}]
  % in
  \begin{scope}[shift={(0,0)}]
    \pelotaLugar{0.7,0.7}
    \fill[white] (0,0) rectangle (1.4,0.62); \draw (0,1.1) -- (0,0) -- (1.4,0) -- (1.4,1.1);
    \draw (0,1.1) -- (1.4,1.1);
    \node[etiqueta] at (0.85,-0.6) {in};
  \end{scope}
  % on
  \begin{scope}[shift={(3.1,0)}]
    \cajaLugar{0,0} \pelotaLugar{0.75,1.55}
    \node[etiqueta] at (0.85,-0.6) {on};
  \end{scope}
  % under
  \begin{scope}[shift={(6.1,0)}]
    \mesaLugar{0,0} \pelotaLugar{1.1,0.36}
    \node[etiqueta] at (1.1,-0.6) {under};
  \end{scope}
  % next to
  \begin{scope}[shift={(9.6,0)}]
    \cajaLugar{0,0} \pelotaLugar{2.3,0.36}
    \node[etiqueta] at (1.3,-0.6) {next to};
  \end{scope}
  % behind
  \begin{scope}[shift={(0.3,-3.6)}]
    \pelotaLugar{1.1,1.35} \cajaLugar{0,0}
    \node[etiqueta] at (0.85,-0.6) {behind};
  \end{scope}
  % in front of
  \begin{scope}[shift={(3.5,-3.6)}]
    \cajaLugar{0,0} \pelotaLugar{0.7,0.3}
    \node[etiqueta] at (0.85,-0.6) {in front of};
  \end{scope}
  % between
  \begin{scope}[shift={(7.2,-3.6)}]
    \cajaLugar{0,0} \pelotaLugar{2.5,0.36} \cajaLugar{3.3,0}
    \node[etiqueta] at (2.55,-0.6) {between};
  \end{scope}
\end{tikzpicture}
\end{center}
\newpage
